\subsection{最後の鏡曲線の全正常切断を一次因子で排除する}
\begin{proposition}
鏡共鳴の残余 $J_m=0$ の両枝にはH3対象はない。
従って鏡共鳴面全体が排除され、両top非零mixed4 $(2,2)$ の未知は $L=0$ の非鏡領域だけとなる。
\end{proposition}
\begin{proof}
$s=0$ は排除済みなので、$\rho(\rho^2-1)(\rho^2+1)\ne0$ として
\[
 s=\frac{i(\rho+1)(\rho^2+1)}{2\rho},\qquad
 t=\frac{\rho^2-1}{2},\qquad
 a=-\frac{(\rho^2-1)^2}{4\rho^2}
\]
の枝を扱えばよい。他方は $i\mapsto-i$ で得られる。
\[
 b=-\frac{2i\rho}{\rho^2-1},\qquad
 Q_2=x^2+b(\rho^2+1)(xy+yx)+b^2y^2
       +\frac{2\rho^2}{\rho^2-1}(uv+vu)
\]
と置くと、
\[
 G=(x+by)Q_2
\]
が各AS4 fiberで正常となる。
また $Q_2$ の $x^2$ 係数は1なので非零であり、domain fiberでは $G\ne0$ である。

正常性を全値で保証する根拠を記す。
入力
\[
 \texttt{cubic\_order2\_22\_secant\_mirror\_J\_input.json}
\]
の35本の多項式四次汎関数は、元の全288関係を消し、指定した35語への制限が非零対角行列となる。
従って有理関数体上で一次独立である。
一方、$\rho=2$ の元関係の厳密消去は $\dim B_4=35$ を与える。
関係行列のrankの半連続性からgeneric四次商次元は35以下、35本の汎関数から35以上なので、正確に35である。
これは一標本から正常性を外挿する議論ではない。
35汎関数が $G$ の全四次正常式を消す恒等式により、正常性は有理関数体上で成立する。
四次商の最大rank開集合は全AS4 fiberを含み、そこで余核は局所自由である。
正常式の係数の分母は $\rho^2-1$ の冪だけなので、前の延長補題で全AS4 fiberへ正常性が延びる。

全正常核の次元も全値で調べる。
同じ多項式汎関数から自由正三次32語への正常行列を作り、候補域で単元となる共通行因子だけを除く。
三つの9minorのgcdは、単元因子を除いて
\[
 \rho^{12}(\rho-1)^{11}(\rho+1)^{10}
 (\rho-i)^2(\rho+i)^2\left(\rho^2+\frac13\right)
\]
となる。従って正常核次元は $3\rho^2+1=0$ を除いて1以下となる。
残る点も
\[
 \rho=\frac{i}{\sqrt3},\quad
 s=\frac{\sqrt3+i}{3},\quad t=-\frac23,\quad a=\frac43
\]
で元関係を $\mathbb Q(\sqrt3,i)$ 上で再消去すると、三次商次元20、四次商次元35、正常rank9となる。
この体の埋込みが全符号・共役を覆う。

従って全AS4-domain候補で正常核は上の非零 $G$ が張る一本である。
全正常切断が実際の一次因子 $x+by$ を持ち、商domain性に反する。
前節の鏡支持縮小と合わせて、鏡共鳴面全体を排除した。

\texttt{verify\_cubic\_order2\_22\_secant\_mirror\_J.py}
は全元関係汎関数、generic独立性、正常式、三minorのgcd、数体有限点を再生する。
有理関数体の消去をブラックボックスとして輸入せず、明示多項式入力の恒等式を検証する。
\end{proof}
